% Methodology-focused English thesis abstract.
% The numerical benchmark comparison is maintained in the backup research reports.
Video polyp segmentation (VPS) is an important task in computer-aided
colonoscopy, with recent advances in segmentation foundation models,
particularly the Segment Anything Model (SAM) family, offering promising
capabilities for lesion delineation and temporal propagation. However,
reliable VPS remains challenging due to variations in lesion appearance,
occlusions, and visually similar surrounding tissues. Furthermore, the
effectiveness of prompt-based segmentation depends heavily on accurate
target initialization, while automatically generated bounding-box prompts
from external detectors such as YOLO may be imprecise or ambiguous.

To address these challenges, this thesis proposes a unified automatic VPS
framework that integrates YOLO-based lesion detection, SAM~2-based promptable
segmentation, and the memory-efficient temporal propagation mechanisms of
Lean-SAM2. The proposed framework addresses three complementary challenges
in video segmentation: (1) reliable target initialization, (2)
discriminative spatial representation learning, and (3) target-consistent
temporal memory management. To this end, three modules are introduced.

First, the \textbf{Feature-Guided Prompt Initialization (FGPI)} module
establishes a target-aware initialization strategy by jointly exploiting
YOLO-generated bounding boxes and early visual representations. Within
each detected region, FGPI constructs a set of spatially constrained,
feature-diverse pseudo-target anchors and employs foreground--background
feature contrast to select informative positive and negative point prompts.
These feature-derived points are combined with the original bounding box
and provided to the SAM~2 prompt encoder, enabling a more discriminative
initialization of the target mask.

Second, the \textbf{Region-Aware Multistage Feature Aggregation (RMFA)}
module enhances spatial representation quality by integrating high-resolution
encoder features with deep, memory-conditioned representations. Through
foreground--background-aware channel attention and region-specific feature
discrimination, RMFA models complementary spatial and semantic information
across feature levels. The resulting representations are used to predict
residual corrections to the segmentation logits, refining target boundaries
and suppressing background interference while preserving the pretrained
mask decoder.

Third, the \textbf{Target-Aware Temporal Memory Refinement (TTMR)}
module formulates a target-conditioned memory selection strategy for video
propagation. By combining target-anchor similarity, memory-token relevance,
and temporal diversity, TTMR prioritizes discriminative historical
representations while reducing redundant information under a constrained
memory budget. This mechanism is designed to maintain target identity across
successive frames and mitigate the accumulation of segmentation errors
associated with appearance variation, occlusion, and propagation drift.
